(1) 机械动量 \(\boldsymbol\pi=\mathbf p-q\mathbf A\) 满足 \([\pi_x,\pi_y]=i\hbar qB\)，不依赖规范选择。令 \(s=\operatorname{sgn}(qB)\)，定义 \(a=(\pi_x+i s\pi_y)/\sqrt{2\hbar|qB|}\)，\(a^\dagger=(\pi_x-i s\pi_y)/\sqrt{2\hbar|qB|}\)，可计算得 \([a,a^\dagger]=1\)。将平方 \(\pi_x^2+\pi_y^2\) 重新表达为振子数算符，就避免了直接求解某一特定规范的波函数。

(2) 因为 \(\pi_x^2+\pi_y^2=2\hbar|qB|(a^\dagger a+1/2)\)，所以 \(H=\hbar\omega_c(a^\dagger a+1/2)\)，其中 \(\omega_c=|q|B/m\)。能级为 \(E_n=\hbar\omega_c(n+1/2)\)，\(n=0,1,\ldots\)，相邻能级能量差 \(\hbar\omega_c\)。该结果是轨道 Landau 量子化，不含电子自旋 Zeeman 分裂。

(3) 磁长度定义为 \(\ell_B=\sqrt{\hbar/(|q|B)}\)。在 Landau 规范中，每个允许的沿边界波数对应一个导引中心，中心间隔的相空间计数给每个 Landau 级态数 \(N_\phi\simeq A/(2\pi\ell_B^2)=|q|BA/(2\pi\hbar)=|\Phi_B|/\Phi_0\)，其中 \(\Phi_0=h/|q|\)。边缘和有限体积修正不计入这个热力学极限中的体态简并数。

(4) 如果填满前 \(\nu\) 个自旋不简并 Landau 级，则 \(N=\nu N_\phi\simeq\nu |q|BA/h\)，面密度为 \(n_{2D}=\nu|q|B/h\)。固定 \(B\) 时态间距与总磁通成正比；当 \(B\to0\)，单个能级简并数趋零，同时能级间隔 \(\hbar|q|B/m\to0\)，大量能级密集起来重新形成无磁场的二维连续能谱，不能只固定一个 \(n\) 而求零磁场全部态。
